% Generated draft: requires mathematical and source review.
\documentclass[11pt]{article}
\usepackage[margin=25mm]{geometry}
\usepackage{fontspec}
\setmainfont{Latin Modern Roman}
\usepackage{amsmath,amssymb,mathtools,unicode-math}
\setmathfont{Latin Modern Math}
\usepackage{enumitem,needspace,xurl,hyperref}
\setlength{\parindent}{0pt}
\setlength{\parskip}{5pt}
\setlength{\emergencystretch}{3em}
\newcommand{\N}{\mathbb{N}}
\newcommand{\Z}{\mathbb{Z}}
\newcommand{\Q}{\mathbb{Q}}
\newcommand{\R}{\mathbb{R}}
\newcommand{\C}{\mathbb{C}}
\newcommand{\F}{\mathbb{F}}
\newcommand{\A}{\mathbb{A}}
\newcommand{\PP}{\mathbb{P}}
\newcommand{\E}{\mathbb{E}}
\newcommand{\eps}{\varepsilon}
\newcommand{\dd}{\,\mathrm{d}}
\begin{document}
\section*{UnsolvedMath ID: 700919}
\textbf{Source problem code:} catalogue entry
\subsection{Definitions and mathematical statement}
Let \(G\) be an arbitrary finite group and let \(p\) be an arbitrary prime. Write \(\F_p=\Z/p\Z\); the source notation \(\Z_p\) means this residue field. For a ring \(R\) with identity, let
\[
U(R)=\{u\in R:\text{there exists }v\in R\text{ with }uv=vu=1\}
\]
be its group of units. The group ring \(RG\) consists of sums \(\sum_{g\in G}a_g g\), with \(a_g\in R\), with coefficientwise addition and multiplication induced by the multiplication in \(G\). Coefficient reduction defines a ring homomorphism
\[
\pi_{G,p}:\Z G\longrightarrow\F_pG,
\qquad
\sum_{g\in G}a_g g\longmapsto\sum_{g\in G}(a_g\bmod p)g,
\]
and hence a group homomorphism
\[
\rho_{G,p}:U(\Z G)\longrightarrow U(\F_pG),
\qquad \rho_{G,p}(u)=\pi_{G,p}(u).
\]
For every such \(G\) and \(p\), determine the kernel \(K_{G,p}=\ker\rho_{G,p}\), determine the image \(I_{G,p}=\operatorname{im}\rho_{G,p}\), and give an explicit transversal for the cosets of \(K_{G,p}\) in \(U(\Z G)\). Here a transversal means a subset \(T_{G,p}\subseteq U(\Z G)\) containing exactly one representative of each coset; equivalently, the restriction \(\rho_{G,p}|_{T_{G,p}}:T_{G,p}\to I_{G,p}\) must be a bijection. The kernel is normal, so left and right cosets coincide. The source requests an explicit description but specifies no particular representation of the answer or computational complexity requirement.
\subsection{Short English statement}
For every finite group \(G\) and prime \(p\), describe which integral group-ring units reduce to the identity modulo \(p\), which units over \(\F_p\) arise by reduction, and an explicit choice of one integral unit for each element of the image.
\subsection{Sources}
\begin{itemize}
\item Local UnsolvedMath registry, permanent ID \(700919\), source problem code \texttt{catalogue entry}, titled \emph{Kourovka Notebook Problem 13.1}. The supplied record labels the problem open; this label is not an independently verified assessment.
\item \emph{The Kourovka Notebook (Unsolved Problems in Group Theory)}, original problem 13.1, 13th issue (1995), proposed by R. Zh. Aleev. The supplied background identifies the 21st issue edition (2026), \texttt{arXiv:1401.0300v48}. The supplied collection description names E. I. Khukhro and V. D. Mazurov as editors of the collection covering Issues 1--21 (1965--2026). These bibliographic details are supplied attribution and have not been independently checked.
\end{itemize}
\end{document}
